% Zeszyt ćwiczeń dla zespołów — warsztat „Autonomiczne pojazdy sterowane AI”
% Treść ćwiczeń: cw00.tex … cw08.tex (wygenerowane z ../*.md skryptem _md_do_latex.py, można edytować ręcznie).
% Kompilacja: pdflatex zeszyt_cwiczen.tex (2×, dla spisu treści); działa też xelatex / lualatex / tectonic
\documentclass[11pt,a4paper,parskip=half]{scrartcl}

% ---------- czcionki: działa z pdfLaTeX (klasycznie: T1 + inputenc) oraz z XeLaTeX/LuaLaTeX ----------
\newif\ifpagella
\pagellafalse   % \pagellatrue → TeX Gyre Pagella + Heros (odpowiedniki Palatino i Helvetiki); domyślnie Latin Modern
\usepackage{iftex}
\ifPDFTeX
  \usepackage[T1]{fontenc}
  \usepackage[utf8]{inputenc}
  \usepackage{lmodern}
  \ifpagella\usepackage{tgpagella}\usepackage[scale=0.92]{tgheros}\fi
  \usepackage{newunicodechar}
  \newunicodechar{→}{\ensuremath{\rightarrow}}
  \newunicodechar{−}{\ensuremath{-}}
  \newunicodechar{▶}{\ensuremath{\blacktriangleright}}
\else
  \usepackage{fontspec}   % bez dodatkowych ustawień: Latin Modern
  \ifpagella
    \setmainfont{texgyrepagella}[Extension=.otf,UprightFont=*-regular,BoldFont=*-bold,ItalicFont=*-italic,BoldItalicFont=*-bolditalic]
    \setsansfont{texgyreheros}[Extension=.otf,UprightFont=*-regular,BoldFont=*-bold,ItalicFont=*-italic,BoldItalicFont=*-bolditalic,Scale=0.92]
  \fi
\fi
\usepackage[polish]{babel}
\usepackage[a4paper,margin=2cm,top=2.3cm,bottom=2.5cm]{geometry}
\usepackage{xcolor,tikz,longtable,booktabs,array,calc,enumitem,pifont,amssymb}
\usepackage[most]{tcolorbox}
\usepackage[hidelinks]{hyperref}
\usepackage{scrlayer-scrpage}

% ---------- kolory (jak w scenariuszu prowadzącego) ----------
\definecolor{morski}{HTML}{1F5F7A}
\definecolor{tlo}{HTML}{F3F7F9}
\definecolor{zolty}{HTML}{E1A100}

% ---------- nagłówki, stopka ----------
\setcounter{secnumdepth}{0}                       % numery są w tytułach („Ćwiczenie 3: …”)
\setcounter{tocdepth}{1}                          % spis treści: tylko ćwiczenia i główne części
\addtokomafont{disposition}{\color{morski}}
\setkomafont{pageheadfoot}{\normalfont\sffamily\footnotesize\color{gray}}
\pagestyle{scrheadings}
\clearpairofpagestyles
\ihead{Autonomiczne pojazdy sterowane AI — zeszyt ćwiczeń}
\ohead{\pagemark}
\ifoot{Zespół nr \rule{1.5cm}{0.4pt}}
\RedeclareSectionCommand[beforeskip=0pt,afterskip=0.8\baselineskip]{section}

% ---------- elementy używane przez treść z pandoca ----------
\setlength{\emergencystretch}{3em}                % długie nazwy funkcji w \texttt nie wystają za margines
\providecommand{\tightlist}{\setlength{\itemsep}{0pt}\setlength{\parskip}{0pt}}
\renewcommand{\arraystretch}{1.75}                % wyższe wiersze tabel — miejsce na pisanie odręczne
\setlength{\LTpre}{4pt}\setlength{\LTpost}{6pt}
\newcommand{\linia}[1]{\rule[-2pt]{#1em}{0.4pt}}  % linia do wpisania odpowiedzi
\renewenvironment{quote}%
  {\begin{tcolorbox}[enhanced,breakable,colback=zolty!8,colframe=zolty!80!black,boxrule=0.6pt,left=6pt,right=6pt]}%
  {\end{tcolorbox}}
\newcommand{\szkicmapy}{%
  \begin{tcolorbox}[colback=white,colframe=gray!60,boxrule=0.5pt,height=7cm,title={\sffamily\small Szkic naszej najlepszej mapy},
    coltitle=black,colbacktitle=gray!15]\end{tcolorbox}}

\begin{document}

% =====================================================================================
\begin{titlepage}
\begin{tikzpicture}[remember picture,overlay]
  \fill[morski] (current page.north west) rectangle ([yshift=-8.5cm]current page.north east);
  \node[anchor=north west,text=white,font=\sffamily\bfseries\fontsize{30}{36}\selectfont,align=left]
    at ([xshift=2cm,yshift=-2.4cm]current page.north west) {Autonomiczne pojazdy\\sterowane AI};
  \node[anchor=north west,text=white!85!morski,font=\sffamily\Large,align=left]
    at ([xshift=2cm,yshift=-5.6cm]current page.north west) {Zeszyt ćwiczeń dla zespołów\\[2pt]\normalsize warsztat 180 min};
\end{tikzpicture}
\vspace*{8.5cm}

{\sffamily\large
\renewcommand{\arraystretch}{2.2}
\begin{tabular}{@{}ll@{}}
Zespół nr & \rule{2cm}{0.4pt}\\
Skład zespołu & \rule{9cm}{0.4pt}\\
              & \rule{9cm}{0.4pt}\\
Robot (Jetbot) nr & \rule{2cm}{0.4pt}\\
Tor & \rule{2cm}{0.4pt}\\
\end{tabular}}

\vfill
{\sffamily\small
\begin{tcolorbox}[colback=tlo,colframe=morski,boxrule=0.6pt]
\textbf{Jak korzystać z zeszytu.} Każde ćwiczenie ma swój notatnik Jupyter o tym samym numerze. Wyniki wpisujcie
od razu do tabel, bo przydadzą się w podsumowaniu (ćwiczenie 8). Zmieniajcie tylko linie oznaczone
\texttt{\# ZMIEŃ TUTAJ}. Gdy robot jedzie, pilot jest zawsze przy torze. Hamulec: komórka \texttt{stop()}.
\end{tcolorbox}}
\end{titlepage}

\tableofcontents
\clearpage

% Wygenerowano z ../00_wprowadzenie.md skryptem _md_do_latex.py (można edytować ręcznie).
\section{Jak pracujemy}\label{jak-pracujemy}

Pracujecie w zespołach 2--3-osobowych. Każdy zespół ma \textbf{jednego Jetbota}, \textbf{stanowisko komputerowe} i \textbf{ten zeszyt ćwiczeń}. Wyniki wpisujecie bezpośrednio do tabel w instrukcjach. Przydadzą się w podsumowaniu.

Podział pracy ustalacie sami. Pilnujcie tylko dwóch zasad:

\begin{itemize}
\tightlist
\item
  \textbf{przy klawiaturze zmieniajcie się co ćwiczenie}, żeby każdy coś uruchomił i zmienił w notatniku,
\item
  \textbf{gdy robot jedzie, jedna osoba stoi przy torze} (w instrukcjach nazywamy ją \textbf{pilotem}). Ustawia robota na starcie, a w razie zagrożenia łapie go albo woła „STOP!''.
\end{itemize}

\subsection{Ćwiczenia}\label{ux107wiczenia}

\par\noindent\begin{minipage}{\linewidth}
\begin{tabular}{@{}
  >{\raggedright\arraybackslash}p{(\linewidth - 6\tabcolsep) * \real{0.1053}}
  >{\raggedright\arraybackslash}p{(\linewidth - 6\tabcolsep) * \real{0.5684}}
  >{\raggedright\arraybackslash}p{(\linewidth - 6\tabcolsep) * \real{0.2211}}
  >{\raggedright\arraybackslash}p{(\linewidth - 6\tabcolsep) * \real{0.1053}}@{}}
\toprule\noalign{}
\begin{minipage}[b]{\linewidth}\raggedright
Nr
\end{minipage} & \begin{minipage}[b]{\linewidth}\raggedright
Ćwiczenie
\end{minipage} & \begin{minipage}[b]{\linewidth}\raggedright
Notatnik
\end{minipage} & \begin{minipage}[b]{\linewidth}\raggedright
Czas
\end{minipage} \\
\midrule\noalign{}
0 & Pierwsze kroki z notatnikiem i robotem & \texttt{00\_start.ipynb} & 5 min \\
1 & LiDAR: jak robot mierzy odległość & \texttt{01\_lidar.ipynb} & 12 min \\
2 & Kamera: jak zapisany jest obraz & \texttt{02\_kamera.ipynb} & 8 min \\
3 & Odometria: dlaczego robot się gubi & \texttt{03\_odometria.ipynb} & 15 min \\
4 & SLAM: budujemy mapę & \texttt{04\_slam.ipynb} & 40 min \\
& \emph{przerwa} & & \\
5 & Neuron: hamuj czy jedź? & \texttt{05\_perceptron.ipynb} & 10 min \\
6 & Gotowa sieć neuronowa szuka obiektów: jak zapytać sieć & \texttt{06\_detektor.ipynb} & 20 min \\
7 & Autonomiczna jazda na mapie i poszukiwanie kaczki & \texttt{07\_nawigacja.ipynb} & 35 min \\
8 & Karta wyników zespołu & (papier) & 20 min \\
\bottomrule\noalign{}
\end{tabular}
\end{minipage}\par\medskip

\section{Ćwiczenie 0: pierwsze kroki (5 min)}\label{ux107wiczenie-0-pierwsze-kroki-5-min}

\textbf{Cel:} nauczyć się uruchamiać komórki notatnika i sprawdzić, czy robot jest połączony.

\subsection{Jak działa notatnik Jupyter}\label{jak-dziaux142a-notatnik-jupyter}

\begin{itemize}
\tightlist
\item
  Notatnik składa się z \textbf{komórek}. Komórki z tekstem to opis, a komórki z kodem to polecenia dla komputera.
\item
  Komórkę uruchamiasz, gdy klikniesz ją i naciśniesz \textbf{Shift + Enter} (albo przycisk ▶).
\item
  Gdy komórka pracuje, obok niej widać \texttt{{[}*{]}}. Po zakończeniu pojawia się numer, np. \texttt{{[}3{]}}.
\item
  Komórki uruchamiaj \textbf{po kolei, od góry}.
\item
  Wolno zmieniać tylko wartości w liniach oznaczonych \textbf{\texttt{\#\ ZMIEŃ\ TUTAJ}}. Resztę kodu można czytać, ale nie trzeba go rozumieć.
\item
  Gdy coś „zawiśnie'', wybierz z menu \emph{Kernel → Restart} i uruchom komórki od początku.
\item
  \textbf{Ostatnia komórka każdego notatnika to \texttt{koniec()}.} Uruchom ją, zanim otworzysz następny notatnik: zatrzymuje robota i odłącza notatnik. Robot działa dalej, a następny notatnik połączy się z nim w pierwszej komórce. Niczego nie trzeba restartować ani zamykać.
\end{itemize}

\subsection{Kroki}\label{kroki}

\begin{enumerate}
\def\labelenumi{\arabic{enumi}.}
\tightlist
\item
  Otwórz \texttt{00\_start.ipynb}. Uruchom pierwszą komórkę (połączenie z robotem). Ma pojawić się komunikat \emph{„Połączono z robotem''}. \textbf{Sprawdź, czy numer w komunikacie, np. \texttt{{[}robot 3{]}}, zgadza się z naklejką na Waszym robocie.} Jeśli nie, zawołaj prowadzącego.
\item
  Uruchom komórkę \texttt{status()}. Przy napędzie, LiDAR-ze i kamerze powinien być znak \ding{51}.
\item
  W komórce z krótkim przejazdem zmień czas jazdy (\texttt{czas\ =\ 1.0\ \ \#\ ZMIEŃ\ TUTAJ}) na \texttt{2.0} i uruchom ją. \textbf{Pilot} trzyma rękę przy robocie.
\item
  Uruchom komórkę z \texttt{stop()}. Zapamiętaj ją, bo to \textbf{hamulec awaryjny}, który działa w każdym notatniku.
\end{enumerate}

\begin{quote}
\textbf{Bezpieczeństwo.} Robot jest mały, ale potrafi spaść ze stołu albo wjechać pod nogi. Gdy robot jedzie, pilot jest zawsze przy torze. W razie problemu: komórka \texttt{stop()} albo podnieś robota.
\end{quote}

\clearpage% Wygenerowano z ../01_lidar.md skryptem _md_do_latex.py (można edytować ręcznie).
\section{Ćwiczenie 1: LiDAR, czyli jak robot mierzy odległość (12 min)}\label{ux107wiczenie-1-lidar-czyli-jak-robot-mierzy-odlegux142oux15bux107-12-min}

\textbf{Notatnik:} \texttt{01\_lidar.ipynb}

\textbf{Potrzebne:} miarka, karton (ok. 30 × 30 cm), zestaw materiałów: biała kartka, czarna matowa tkanina, szyba lub plexi, lusterko, metalowa puszka, kątomierz i linijka (do ustawiania kątów)

\subsection{Cel}\label{cel}

Sprawdzić, jak dokładnie LiDAR mierzy odległość i w jakich sytuacjach się myli.

\subsection{Trochę teorii}\label{trochux119-teorii}

LiDAR (\emph{Light Detection and Ranging}) wysyła krótkie impulsy światła laserowego i mierzy, po jakim czasie wracają po odbiciu od przeszkody. Ponieważ światło porusza się ze stałą prędkością, z czasu powrotu można policzyć odległość. Głowica LiDAR-u obraca się kilka razy na sekundę. W ten sposób robot dostaje \textbf{kilkaset pomiarów odległości dookoła siebie}. Taki zestaw punktów to \textbf{skan} (chmura punktów).

\begin{itemize}
\tightlist
\item
  \textbf{LiDAR 2D} (taki jak na Jetbocie) mierzy w jednej płaszczyźnie, na wysokości czujnika. Nie widzi więc przedmiotów niższych lub wyższych od tej płaszczyzny.
\item
  \textbf{LiDAR 3D} ma wiele wiązek ustawionych pod różnymi kątami i widzi otoczenie w przestrzeni. Takie czujniki stosuje się w samochodach autonomicznych.
\end{itemize}

\subsection{Kroki}\label{kroki}

\subsubsection{Krok 1: pierwszy skan (2 min)}\label{krok-1-pierwszy-skan-2-min}

Uruchom komórki \texttt{polacz()} i \texttt{rysuj\_skan()}. Widzisz otoczenie robota z góry. Robot to czerwony trójkąt, a jego przód jest skierowany w górę wykresu.

\begin{itemize}
\tightlist
\item
  Gdzie na wykresie są ściany? Gdzie przeszkody na torze?
\item
  Jedna osoba staje przed robotem, a druga ponownie uruchamia \texttt{rysuj\_skan()}. Jak na skanie wyglądają nogi?
\end{itemize}

\subsubsection{Krok 2: dokładność pomiaru (4 min)}\label{krok-2-dokux142adnoux15bux107-pomiaru-4-min}

Ustaw karton \textbf{na wprost robota}, prostopadle do niego. Odległość mierz miarką od \textbf{środka LiDAR-u} (obrotowej głowicy) do kartonu. Dla każdej odległości uruchom \texttt{seria\_pomiarow(kat=0)}. Komórka wykona 10 pomiarów i poda średnią oraz rozrzut.

\par\noindent\begin{minipage}{\linewidth}
\begin{tabular}{@{}
  >{\raggedright\arraybackslash}p{(\linewidth - 6\tabcolsep) * \real{0.2366}}
  >{\raggedright\arraybackslash}p{(\linewidth - 6\tabcolsep) * \real{0.2258}}
  >{\raggedright\arraybackslash}p{(\linewidth - 6\tabcolsep) * \real{0.2581}}
  >{\raggedright\arraybackslash}p{(\linewidth - 6\tabcolsep) * \real{0.2796}}@{}}
\toprule\noalign{}
\begin{minipage}[b]{\linewidth}\raggedright
Odległość z miarki {[}m{]}
\end{minipage} & \begin{minipage}[b]{\linewidth}\raggedright
Średnia z LiDAR-u {[}m{]}
\end{minipage} & \begin{minipage}[b]{\linewidth}\raggedright
Rozrzut (odch. std.) {[}m{]}
\end{minipage} & \begin{minipage}[b]{\linewidth}\raggedright
Błąd = LiDAR − miarka {[}cm{]}
\end{minipage} \\
\midrule\noalign{}
0,30 & & & \\
0,60 & & & \\
1,00 & & & \\
1,50 & & & \\
\bottomrule\noalign{}
\end{tabular}
\end{minipage}\par\medskip

\subsubsection{Krok 3: różne materiały (3 min)}\label{krok-3-ruxf3ux17cne-materiaux142y-3-min}

Na kartonie ustawionym w odległości \textbf{0,60 m} umieszczaj kolejno różne materiały. \textbf{Lusterko} zmierz dwa razy: ustawione na wprost robota i obrócone o ok. 45°. Przy lusterku pod kątem uruchom też \texttt{rysuj\_skan()}: gdzie jest punkt?

\par\noindent\begin{minipage}{\linewidth}
\begin{tabular}{@{}
  >{\raggedright\arraybackslash}p{(\linewidth - 6\tabcolsep) * \real{0.2222}}
  >{\raggedright\arraybackslash}p{(\linewidth - 6\tabcolsep) * \real{0.2121}}
  >{\raggedright\arraybackslash}p{(\linewidth - 6\tabcolsep) * \real{0.3030}}
  >{\raggedright\arraybackslash}p{(\linewidth - 6\tabcolsep) * \real{0.2626}}@{}}
\toprule\noalign{}
\begin{minipage}[b]{\linewidth}\raggedright
Materiał
\end{minipage} & \begin{minipage}[b]{\linewidth}\raggedright
Średnia z LiDAR-u {[}m{]}
\end{minipage} & \begin{minipage}[b]{\linewidth}\raggedright
Ile pomiarów nieudanych (z 10)?
\end{minipage} & \begin{minipage}[b]{\linewidth}\raggedright
Czy wynik jest wiarygodny?
\end{minipage} \\
\midrule\noalign{}
karton (odniesienie) & & & \\
biała kartka & & & \\
czarna matowa tkanina & & & \\
szyba / plexi & & & \\
lusterko na wprost & & & \\
lusterko pod kątem 45° & & & \\
metalowa puszka & & & \\
\bottomrule\noalign{}
\end{tabular}
\end{minipage}\par\medskip

\subsubsection{Krok 4: przeszkoda pod kątem (2 min)}\label{krok-4-przeszkoda-pod-kux105tem-2-min}

Karton w odległości \textbf{0,60 m}, obracany względem robota. Kąt ustawcie kątomierzem: linijkę przyłóżcie do przedniej ściany kartonu, kątomierz do linijki (0° = karton ustawiony prostopadle do robota).

\par\noindent\begin{minipage}{\linewidth}
\begin{tabular}{@{}
  >{\raggedright\arraybackslash}p{(\linewidth - 4\tabcolsep) * \real{0.3433}}
  >{\raggedright\arraybackslash}p{(\linewidth - 4\tabcolsep) * \real{0.3134}}
  >{\raggedright\arraybackslash}p{(\linewidth - 4\tabcolsep) * \real{0.3433}}@{}}
\toprule\noalign{}
\begin{minipage}[b]{\linewidth}\raggedright
Kąt ustawienia kartonu
\end{minipage} & \begin{minipage}[b]{\linewidth}\raggedright
Średnia z LiDAR-u {[}m{]}
\end{minipage} & \begin{minipage}[b]{\linewidth}\raggedright
Pomiary nieudane (z 10)
\end{minipage} \\
\midrule\noalign{}
0° (prostopadle) & & \\
30° & & \\
60° & & \\
prawie równolegle (80°) & & \\
\bottomrule\noalign{}
\end{tabular}
\end{minipage}\par\medskip

\subsection{Pytania do przemyślenia}\label{pytania-do-przemyux15blenia}

\begin{enumerate}
\def\labelenumi{\arabic{enumi}.}
\tightlist
\item
  Czy błąd pomiaru rośnie z odległością? O ile centymetrów może się mylić LiDAR?
\item
  Przy których materiałach LiDAR zawodzi? Dlaczego? \emph{(Podpowiedź: co dzieje się ze światłem przy szybie, czarnej tkaninie i lustrze? Dlaczego lustro na wprost i pod kątem dało inne wyniki?)}
\item
  Robot autonomiczny jedzie korytarzem z przeszklonymi drzwiami. Jakie zagrożenie widzisz?
\item
  Stolik ma blat na wysokości 70 cm i cienkie nogi. Co „zobaczy'' LiDAR 2D zamontowany 15 cm nad podłogą?
\end{enumerate}

\begin{quote}
\textbf{Jeśli coś nie działa.} „Brak poprawnych pomiarów'' oznacza, że czujnik nie dostał odbicia. To też wynik, więc wpisz go do tabeli. Jeśli żaden skan się nie pojawia, uruchom \texttt{status()} i zawołaj prowadzącego.
\end{quote}

\clearpage% Wygenerowano z ../02_kamera.md skryptem _md_do_latex.py (można edytować ręcznie).
\section{Ćwiczenie 2: kamera, czyli jak zapisany jest obraz (8 min)}\label{ux107wiczenie-2-kamera-czyli-jak-zapisany-jest-obraz-8-min}

\textbf{Notatnik:} \texttt{02\_kamera.ipynb}

\textbf{Potrzebne:} kolorowe przedmioty przed kamerą robota (np. czerwona gumowa kaczka, zielona i niebieska kartka, kartka z napisem drukowanymi literami)

\subsection{Cel}\label{cel}

Zrozumieć, z czego składa się obraz cyfrowy i dlaczego sposób jego zapisu ma znaczenie dla robota.

\subsection{Trochę teorii}\label{trochux119-teorii}

\begin{itemize}
\tightlist
\item
  Obraz cyfrowy to prostokątna siatka \textbf{pikseli}. \textbf{Rozdzielczość} to liczba pikseli w poziomie × w pionie, np. 640 × 480.
\item
  W modelu \textbf{RGB} kolor każdego piksela jest zapisany jako trzy liczby od 0 do 255: ilość światła \textbf{czerwonego (R)}, \textbf{zielonego (G)} i \textbf{niebieskiego (B)}. (0, 0, 0) to czerń, a (255, 255, 255) to biel.
\item
  Jeden piksel zajmuje 3 bajty. Obraz 640 × 480 zajmuje więc 640 · 480 · 3 = \textbf{921 600 bajtów (900 kB)}. Kamera wysyła ok. 30 takich obrazów na sekundę.
\item
  \textbf{Kompresja bezstratna} (np. \textbf{PNG}) zmniejsza plik bez utraty informacji: po odczycie obraz jest identyczny.
\item
  \textbf{Kompresja stratna} (np. \textbf{JPEG}) usuwa szczegóły, których oko prawie nie zauważa. Pliki są dużo mniejsze, ale obraz już nie jest identyczny z oryginałem.
\end{itemize}

\subsection{Kroki}\label{kroki}

\subsubsection{Krok 1: zdjęcie i rozdzielczość (1 min)}\label{krok-1-zdjux119cie-i-rozdzielczoux15bux107-1-min}

Ustaw przedmioty przed kamerą i uruchom komórkę ze zdjęciem.

\begin{itemize}
\tightlist
\item
  Rozdzielczość obrazu: \linia{2.5} × \linia{2.5} pikseli; liczba pikseli: \linia{5.0}
\item
  Rozmiar bez kompresji: \linia{2.5} kB
\end{itemize}

\subsubsection{Krok 2: piksele i kolory (2 min)}\label{krok-2-piksele-i-kolory-2-min}

W komórce \texttt{piksele(img,\ x,\ y)} zmień \texttt{x} i \texttt{y} tak, żeby oglądać fragment \textbf{kaczki}, a potem \textbf{białej ściany lub kartki}. Wartości odczytasz z powiększenia.

\par\noindent\begin{minipage}{\linewidth}
\begin{tabular}{@{}
  >{\raggedright\arraybackslash}p{(\linewidth - 6\tabcolsep) * \real{0.3043}}
  >{\raggedright\arraybackslash}p{(\linewidth - 6\tabcolsep) * \real{0.2319}}
  >{\raggedright\arraybackslash}p{(\linewidth - 6\tabcolsep) * \real{0.2319}}
  >{\raggedright\arraybackslash}p{(\linewidth - 6\tabcolsep) * \real{0.2319}}@{}}
\toprule\noalign{}
\begin{minipage}[b]{\linewidth}\raggedright
Miejsce
\end{minipage} & \begin{minipage}[b]{\linewidth}\raggedright
R
\end{minipage} & \begin{minipage}[b]{\linewidth}\raggedright
G
\end{minipage} & \begin{minipage}[b]{\linewidth}\raggedright
B
\end{minipage} \\
\midrule\noalign{}
kaczka (czerwona) & & & \\
biała ściana / kartka & & & \\
ciemny fragment & & & \\
\bottomrule\noalign{}
\end{tabular}
\end{minipage}\par\medskip

Uruchom \texttt{kanaly\_rgb(img)}. W którym kanale kaczka jest najjaśniejsza? \linia{5.0}

\subsubsection{Krok 3: kompresja (3 min)}\label{krok-3-kompresja-3-min}

Uruchom \texttt{porownaj\_kompresje(img)}. Potem uruchom ją jeszcze raz, podając \texttt{wycinek}, czyli powiększony fragment z krawędzią kaczki.

\par\noindent\begin{minipage}{\linewidth}
\begin{tabular}{@{}
  >{\raggedright\arraybackslash}p{(\linewidth - 6\tabcolsep) * \real{0.1772}}
  >{\raggedright\arraybackslash}p{(\linewidth - 6\tabcolsep) * \real{0.2025}}
  >{\raggedright\arraybackslash}p{(\linewidth - 6\tabcolsep) * \real{0.2405}}
  >{\raggedright\arraybackslash}p{(\linewidth - 6\tabcolsep) * \real{0.3797}}@{}}
\toprule\noalign{}
\begin{minipage}[b]{\linewidth}\raggedright
Format
\end{minipage} & \begin{minipage}[b]{\linewidth}\raggedright
Rozmiar {[}kB{]}
\end{minipage} & \begin{minipage}[b]{\linewidth}\raggedright
Średni błąd piksela
\end{minipage} & \begin{minipage}[b]{\linewidth}\raggedright
Czy widać różnicę na powiększeniu?
\end{minipage} \\
\midrule\noalign{}
bez kompresji & & 0 & --- \\
PNG & & & \\
JPEG jakość 95 & & & \\
JPEG jakość 50 & & & \\
JPEG jakość 10 & & & \\
\bottomrule\noalign{}
\end{tabular}
\end{minipage}\par\medskip

\subsubsection{Krok 4: mniejsza rozdzielczość (2 min)}\label{krok-4-mniejsza-rozdzielczoux15bux107-2-min}

W komórce \texttt{zmien\_rozdzielczosc} sprawdź szerokości 320, 160, 80 i 40 pikseli.

\begin{itemize}
\tightlist
\item
  Przy jakiej szerokości napis na kartce przestaje być czytelny? \linia{2.5} px A kaczka przestaje być rozpoznawalna? \linia{2.5} px
\end{itemize}

\subsection{Pytania do przemyślenia}\label{pytania-do-przemyux15blenia}

\begin{enumerate}
\def\labelenumi{\arabic{enumi}.}
\tightlist
\item
  Robot wysyła obraz z kamery przez Wi-Fi do komputera stacji, na którym działa sieć neuronowa. Który format wybierzesz do przesyłania na żywo, a który do zapisania zdjęć dokumentujących wypadek? Dlaczego?
\item
  Ile danych na sekundę wysyłałaby kamera 640 × 480 nadająca 30 klatek/s bez kompresji?
\item
  Mniejszy obraz to szybsze przetwarzanie. Co tracimy? \emph{(Wrócimy do tego w ćwiczeniu 6.)}
\end{enumerate}

\clearpage% Wygenerowano z ../03_odometria.md skryptem _md_do_latex.py (można edytować ręcznie).
\section{Ćwiczenie 3: odometria, czyli dlaczego robot się gubi (15 min)}\label{ux107wiczenie-3-odometria-czyli-dlaczego-robot-siux119-gubi-15-min}

\textbf{Notatnik:} \texttt{03\_odometria.ipynb}

\textbf{Potrzebne:} taśma malarska, miarka, kawałek wolnej podłogi ok. 1,5 × 1,5 m

\subsection{Cel}\label{cel}

Sprawdzić, jak dokładnie robot zna swoje położenie, gdy liczy je tylko z obrotów kół.

\subsection{Trochę teorii}\label{trochux119-teorii}

\textbf{Odometria} to obliczanie położenia na podstawie tego, ile obrotów wykonały koła. Robot „wie'', że koło o obwodzie 20 cm obróciło się 5 razy, więc uznaje, że przejechał 1 m. Tak samo liczy obroty. Problem w tym, że koła się \textbf{ślizgają}, a dwa silniki \textbf{nigdy nie są identyczne}. Każdy mały błąd dokłada się do poprzednich, a robot nie ma jak tego zauważyć.

\subsection{Kroki}\label{kroki}

\subsubsection{Krok 1: oznaczenie startu (2 min)}\label{krok-1-oznaczenie-startu-2-min}

Przyklej na podłodze taśmą \textbf{krzyżyk} (punkt startu) i \textbf{linię kierunku}. Ustaw robota tak, żeby środek między kołami był nad krzyżykiem, a przód był skierowany wzdłuż linii. Przyklej też kawałek taśmy przy przednim brzegu robota, żeby potem zmierzyć kąt.

\subsubsection{Krok 2: przejazd po kwadracie (4 min)}\label{krok-2-przejazd-po-kwadracie-4-min}

Uruchom \texttt{przejazd\_po\_kwadracie(bok=1.0)}. Robot przejedzie kwadrat 1 × 1 m, skręcając w lewo, i zatrzyma się tam, gdzie \textbf{według odometrii} jest start.

Gdy robot stanie, zmierzcie miarką, jak daleko środek robota jest od krzyżyka (w cm). Pomiar obrotu robota względem linii kierunku jest szacunkowy.

\subsubsection{Krok 3: kolejne przejazdy (6 min)}\label{krok-3-kolejne-przejazdy-6-min}

Za każdym razem ustaw robota dokładnie na krzyżyku. W komórkach zmieniasz parametry oznaczone \texttt{\#\ ZMIEŃ\ TUTAJ}.

\par\noindent\begin{minipage}{\linewidth}
\begin{tabular}{@{}
  >{\raggedright\arraybackslash}p{(\linewidth - 12\tabcolsep) * \real{0.0752}}
  >{\raggedright\arraybackslash}p{(\linewidth - 12\tabcolsep) * \real{0.0752}}
  >{\raggedright\arraybackslash}p{(\linewidth - 12\tabcolsep) * \real{0.1053}}
  >{\raggedright\arraybackslash}p{(\linewidth - 12\tabcolsep) * \real{0.1128}}
  >{\raggedright\arraybackslash}p{(\linewidth - 12\tabcolsep) * \real{0.2256}}
  >{\raggedright\arraybackslash}p{(\linewidth - 12\tabcolsep) * \real{0.2256}}
  >{\raggedright\arraybackslash}p{(\linewidth - 12\tabcolsep) * \real{0.1805}}@{}}
\toprule\noalign{}
\begin{minipage}[b]{\linewidth}\raggedright
Przejazd
\end{minipage} & \begin{minipage}[b]{\linewidth}\raggedright
Bok {[}m{]}
\end{minipage} & \begin{minipage}[b]{\linewidth}\raggedright
Prędkość {[}m/s{]}
\end{minipage} & \begin{minipage}[b]{\linewidth}\raggedright
Okrążenia
\end{minipage} & \begin{minipage}[b]{\linewidth}\raggedright
Odległość od startu wg odometrii {[}cm{]}
\end{minipage} & \begin{minipage}[b]{\linewidth}\raggedright
Odległość od startu z miarki {[}cm{]}
\end{minipage} & \begin{minipage}[b]{\linewidth}\raggedright
Obrót względem linii {[}°{]}
\end{minipage} \\
\midrule\noalign{}
1 & 1,0 & 0,15 & 1 & & & \\
2 & 1,0 & 0,15 & 1 (powtórzenie) & & & \\
3 & 1,0 & 0,30 & 1 & & & \\
4 & 1,0 & 0,15 & 2 & & & \\
\bottomrule\noalign{}
\end{tabular}
\end{minipage}\par\medskip

\subsubsection{Krok 4: trasa według robota (1 min)}\label{krok-4-trasa-wedux142ug-robota-1-min}

Uruchom \texttt{rysuj\_trase(...)}. Wykres pokazuje, jak robot \textbf{wyobraża sobie} swoją trasę.

\subsection{Pytania do przemyślenia}\label{pytania-do-przemyux15blenia}

\begin{enumerate}
\def\labelenumi{\arabic{enumi}.}
\tightlist
\item
  Robot „sądzi'', że wrócił do startu. Gdzie jest naprawdę? Czy przy powtórzeniu błąd jest taki sam?
\item
  Jak zmienia się błąd przy większej prędkości? A przy dwóch okrążeniach?
\item
  Gdyby robot przejechał tak 100 m, czy nadal wiedziałby, gdzie jest?
\item
  Czego robot potrzebuje, żeby \textbf{poprawić} swoje położenie? \emph{(Podpowiedź: co widzi LiDAR? Zobacz ćwiczenie 4.)}
\end{enumerate}

\begin{quote}
\textbf{Wniosek do zapamiętania.} Odometria jest dobra na krótkich odcinkach, ale jej błąd \textbf{narasta bez końca}. Dlatego robot musi porównywać to, co widzi, z mapą otoczenia. Na tym polega \textbf{SLAM}.
\end{quote}

\clearpage% Wygenerowano z ../04_slam.md skryptem _md_do_latex.py (można edytować ręcznie).
\section{Ćwiczenie 4: SLAM, czyli budujemy mapę (40 min)}\label{ux107wiczenie-4-slam-czyli-budujemy-mapux119-40-min}

\textbf{Notatnik:} \texttt{04\_slam.ipynb}

\textbf{Potrzebne:} tor z przeszkodami, punkt startowy \textbf{S} Waszego toru oznaczony na podłodze (wspólny dla zespołów jeżdżących na tym torze)

\subsection{Cel}\label{cel}

Zbudować mapę toru metodą SLAM, sprawdzić, jak parametry algorytmu wpływają na jakość mapy, i zapisać najlepszą mapę do części II.

\subsection{Trochę teorii}\label{trochux119-teorii}

\textbf{SLAM} (\emph{Simultaneous Localization and Mapping}) to jednoczesna lokalizacja i budowa mapy. Robot rysuje mapę z kolejnych skanów LiDAR-u. Jednocześnie porównuje każdy nowy skan z tym, co już narysował, i na tej podstawie \textbf{poprawia swoje położenie}, czyli koryguje błąd odometrii z ćwiczenia 3.

Mapa to siatka kratek. Każda kratka może być:

\begin{itemize}
\tightlist
\item
  \textbf{wolna} (biała): tu nie ma przeszkody,
\item
  \textbf{zajęta} (czarna): tu jest przeszkoda,
\item
  \textbf{nieznana} (szara): robot jeszcze tu nie „zajrzał''.
\end{itemize}

\textbf{Zamknięcie pętli} (\emph{loop closure}) zachodzi, gdy robot wraca w miejsce, które już zna. Algorytm zauważa, że „to tu już było'', i prostuje całą mapę.

\textbf{Jak zbudowany jest robot w ROS (Robot Operating System).} Oprogramowanie robota to wiele osobnych programów, czyli \textbf{węzłów}. Jeden obsługuje LiDAR, drugi silniki, trzeci SLAM. Wasz notatnik na stacji też jest węzłem. Węzły wymieniają się wiadomościami przez \textbf{tematy}. Na przykład LiDAR publikuje skany w temacie \texttt{/scan}, a SLAM je stamtąd odbiera i publikuje mapę w temacie \texttt{/map}.

\subsection{Zasady jazdy podczas budowy mapy}\label{zasady-jazdy-podczas-budowy-mapy}

\begin{itemize}
\tightlist
\item
  Jedź \textbf{wolno}, a obracaj się jeszcze wolniej. Szybki obrót „rozmazuje'' mapę.
\item
  Objedź cały tor i \textbf{wróć do miejsc, które już robot widział}, żeby zamknąć pętlę.
\item
  Nie stawajcie wszyscy wokół robota, bo LiDAR narysuje Wasze nogi jako ściany.
\end{itemize}

\subsection{Kroki}\label{kroki}

\subsubsection{Krok 1: ROS od środka (3 min)}\label{krok-1-ros-od-ux15brodka-3-min}

Uruchom \texttt{wezly\_i\_tematy()}. Znajdź na liście węzeł LiDAR-u i temat \texttt{/scan}, a także temat, którym wysyłane są polecenia jazdy (\texttt{/cmd\_vel}).

\subsubsection{Krok 2: pierwsza mapa z parametrami domyślnymi (12 min)}\label{krok-2-pierwsza-mapa-z-parametrami-domyux15blnymi-12-min}

\begin{enumerate}
\def\labelenumi{\arabic{enumi}.}
\tightlist
\item
  \textbf{Ustaw robota na punkcie S}, przodem wzdłuż oznaczonej linii. Punkt S będzie środkiem mapy (x = 0, y = 0), a kierunek linii to oś x.
\item
  Uruchom \texttt{uruchom\_slam()}. Komunikat potwierdzi, że SLAM działa.
\item
  Uruchom \texttt{panel\_sterowania(podglad=True)} i objedź tor przyciskami. Obok przycisków zobaczysz mapę na żywo: robot (fioletowa strzałka), skan LiDAR-u (czerwone punkty) i rosnącą mapę. Obserwuj, jak mapa rośnie.
\item
  Obserwuj zamknięcie pętli: co się dzieje z mapą, gdy robot wraca w znane miejsce?
\item
  Uruchom \texttt{pokaz\_mape()} i zapisz wynik do tabeli. Uruchom \texttt{zapisz\_mape("zespolN\_mapa1")}, wstawiając zamiast N numer zespołu.
\end{enumerate}

\subsubsection{Krok 3: eksperymenty z parametrami (18 min)}\label{krok-3-eksperymenty-z-parametrami-18-min}

Prowadzący przydzieli Waszemu zespołowi \textbf{dwa eksperymenty} z listy poniżej. Przed każdym eksperymentem: robot na punkt S, uruchomić komórkę z nowymi parametrami (SLAM zaczyna od pustej mapy), objechać tor tak samo jak w kroku 2.

\par\noindent\begin{minipage}{\linewidth}
\begin{tabular}{@{}
  >{\raggedright\arraybackslash}p{(\linewidth - 4\tabcolsep) * \real{0.1183}}
  >{\raggedright\arraybackslash}p{(\linewidth - 4\tabcolsep) * \real{0.3118}}
  >{\raggedright\arraybackslash}p{(\linewidth - 4\tabcolsep) * \real{0.5699}}@{}}
\toprule\noalign{}
\begin{minipage}[b]{\linewidth}\raggedright
Eksperyment
\end{minipage} & \begin{minipage}[b]{\linewidth}\raggedright
Co zmieniamy
\end{minipage} & \begin{minipage}[b]{\linewidth}\raggedright
Parametry w notatniku
\end{minipage} \\
\midrule\noalign{}
A & grubsza siatka & \texttt{rozdzielczosc=0.10} \\
B & drobniejsza siatka & \texttt{rozdzielczosc=0.02} \\
C & „krótkowzroczny'' LiDAR & \texttt{zasieg=1.0} \\
D & rzadkie aktualizacje mapy & \texttt{aktualizacja=10.0} \\
E & mało hipotez położenia & \texttt{czastki=5} \\
F & szybka jazda i szybkie obroty & domyślne parametry; w panelu prędkość 0,3 i obrót 2,0 \\
\bottomrule\noalign{}
\end{tabular}
\end{minipage}\par\medskip

\par\noindent\begin{minipage}{\linewidth}
\begin{tabular}{@{}
  >{\raggedright\arraybackslash}p{(\linewidth - 12\tabcolsep) * \real{0.0746}}
  >{\raggedright\arraybackslash}p{(\linewidth - 12\tabcolsep) * \real{0.2388}}
  >{\raggedright\arraybackslash}p{(\linewidth - 12\tabcolsep) * \real{0.1343}}
  >{\raggedright\arraybackslash}p{(\linewidth - 12\tabcolsep) * \real{0.1418}}
  >{\raggedright\arraybackslash}p{(\linewidth - 12\tabcolsep) * \real{0.1716}}
  >{\raggedright\arraybackslash}p{(\linewidth - 12\tabcolsep) * \real{0.1194}}
  >{\raggedright\arraybackslash}p{(\linewidth - 12\tabcolsep) * \real{0.1194}}@{}}
\toprule\noalign{}
\begin{minipage}[b]{\linewidth}\raggedright
Mapa
\end{minipage} & \begin{minipage}[b]{\linewidth}\raggedright
Parametry
\end{minipage} & \begin{minipage}[b]{\linewidth}\raggedright
Czas objazdu {[}min{]}
\end{minipage} & \begin{minipage}[b]{\linewidth}\raggedright
Zbadany obszar {[}m²{]}
\end{minipage} & \begin{minipage}[b]{\linewidth}\raggedright
Czy pętla się zamknęła?
\end{minipage} & \begin{minipage}[b]{\linewidth}\raggedright
Ocena mapy 1--5
\end{minipage} & \begin{minipage}[b]{\linewidth}\raggedright
Co jest nie tak?
\end{minipage} \\
\midrule\noalign{}
mapa1 & domyślne (0,05 / 6,0 / 2,0 / 30) & & & & & \\
& & & & & & \\
& & & & & & \\
\bottomrule\noalign{}
\end{tabular}
\end{minipage}\par\medskip

Kryteria oceny mapy: proste ściany (nie „podwójne''), przeszkody w dobrych miejscach, brak dziur, kształt toru zgodny z rzeczywistością.

\subsubsection{Krok 4: najlepsza mapa (7 min)}\label{krok-4-najlepsza-mapa-7-min}

Ustawcie parametry, które dały najlepszą mapę, zbudujcie ją jeszcze raz starannie i zapiszcie jako \textbf{\texttt{zespolN\_najlepsza}}. Ta mapa posłuży w ćwiczeniu 7. Na koniec uruchom \texttt{zatrzymaj\_slam()}.

\subsection{Pytania do przemyślenia}\label{pytania-do-przemyux15blenia}

\begin{enumerate}
\def\labelenumi{\arabic{enumi}.}
\tightlist
\item
  Co się stało z mapą przy szybkiej jeździe i szybkich obrotach? Dlaczego?
\item
  Drobniejsza siatka daje dokładniejszą mapę. Jaka jest cena? \emph{(Pomyśl o pamięci i czasie obliczeń.)}
\item
  Dlaczego długi, pusty korytarz bez charakterystycznych punktów jest trudny dla SLAM?
\item
  \emph{(Jeśli prowadzący pokazał LiDAR 3D)} Co widzi LiDAR 3D, czego nie widzi 2D?
\end{enumerate}

\begin{quote}
\textbf{Plan awaryjny.} Jeśli mapa wyszła słabo, a czas się kończy, prowadzący ma gotową mapę toru \texttt{tor\_prowadzacy}. Wystarczy jej użyć w ćwiczeniu 7.
\end{quote}

\clearpage% Wygenerowano z ../05_perceptron.md skryptem _md_do_latex.py (można edytować ręcznie).
\section{Ćwiczenie 5: neuron, czyli hamuj czy jedź? (10 min)}\label{ux107wiczenie-5-neuron-czyli-hamuj-czy-jedux17a-10-min}

\textbf{Notatnik:} \texttt{05\_perceptron.ipynb} (nie wymaga robota)

\subsection{Cel}\label{cel}

Zobaczyć, jak pojedynczy sztuczny neuron podejmuje decyzję i co to znaczy, że sieć neuronowa „uczy się''.

\subsection{Trochę teorii}\label{trochux119-teorii}

Sztuczny neuron (\textbf{perceptron}) dostaje liczby, np. odległość od przeszkody. Mnoży je przez \textbf{wagi}, dodaje \textbf{bias} (przesunięcie) i sprawdza znak wyniku:

\begin{quote}
s = waga · odległość + bias\\
s \textgreater{} 0 → \textbf{HAMUJ}, w przeciwnym razie → \textbf{JEDŹ}
\end{quote}

Cała „wiedza'' neuronu to te kilka liczb. \textbf{Uczenie} to automatyczne szukanie takich liczb, przy których neuron odpowiada dobrze.

\textbf{Jak czytać wykresy:} \textbf{tło} pokazuje, co robi neuron (czerwone = HAMUJ, zielone = JEDŹ). \textbf{Kropki} pokazują, co powinien zrobić. \textbf{\ding{54}} oznacza błąd neuronu. Pod wykresem jest przykład obliczenia s.

\subsection{Kroki}\label{kroki}

\subsubsection{Część A: jedno wejście, czyli odległość (3 min)}\label{czux119ux15bux107-a-jedno-wejux15bcie-czyli-odlegux142oux15bux107-3-min}

Waga jest ustawiona na −1. Przesuwaj \textbf{bias} i patrz, jak przesuwa się czarna granica. Ustaw go tak, żeby neuron hamował tylko przy czerwonych kropkach (\ding{51} BRAWO).

\begin{itemize}
\tightlist
\item
  Przy jakim biasie się udało? \linia{2.5} Gdzie jest wtedy granica? \linia{2.5} m
\item
  Zmień wagę na \textbf{+1}. Co robi teraz neuron? Dlaczego waga musi być ujemna?
\end{itemize}

\subsubsection{Część B: dwa wejścia, czyli odległość i prędkość (4 min)}\label{czux119ux15bux107-b-dwa-wejux15bcia-czyli-odlegux142oux15bux107-i-prux119dkoux15bux107-4-min}

Przy szybkiej jeździe trzeba hamować wcześniej. Waga odległości zostaje −1. Dobierz \textbf{wagę prędkości} i \textbf{bias} tak, żeby wszystkie 60 sytuacji było dobrze.

\par\noindent\begin{minipage}{\linewidth}
\begin{tabular}{@{}
  >{\raggedright\arraybackslash}p{(\linewidth - 6\tabcolsep) * \real{0.4146}}
  >{\raggedright\arraybackslash}p{(\linewidth - 6\tabcolsep) * \real{0.1951}}
  >{\raggedright\arraybackslash}p{(\linewidth - 6\tabcolsep) * \real{0.1951}}
  >{\raggedright\arraybackslash}p{(\linewidth - 6\tabcolsep) * \real{0.1951}}@{}}
\toprule\noalign{}
\begin{minipage}[b]{\linewidth}\raggedright
\end{minipage} & \begin{minipage}[b]{\linewidth}\raggedright
waga prędkości
\end{minipage} & \begin{minipage}[b]{\linewidth}\raggedright
bias
\end{minipage} & \begin{minipage}[b]{\linewidth}\raggedright
Poprawne (z 60)
\end{minipage} \\
\midrule\noalign{}
Nasze ustawienie & & & \\
Komputer (część C, po podzieleniu) & & & \\
\bottomrule\noalign{}
\end{tabular}
\end{minipage}\par\medskip

\subsubsection{Część C: automatyczne uczenie (3 min)}\label{czux119ux15bux107-c-automatyczne-uczenie-3-min}

Uruchom komórkę. Komputer zaczyna od wag równych zero i po każdej pomyłce lekko je poprawia. Na wykresach widać, jak granica przesuwa się po kolejnych \textbf{epokach} (epoka to jedno przejście przez wszystkie przykłady). Wpisz wynik do tabeli powyżej.

\subsection{Pytania do przemyślenia}\label{pytania-do-przemyux15blenia}

\begin{enumerate}
\def\labelenumi{\arabic{enumi}.}
\tightlist
\item
  Czy komputer znalazł te same liczby co Wy? Czy to źle, że są inne? \emph{(Porównaj, gdzie przebiega granica.)}
\item
  Wasz neuron ma \textbf{3} liczby do dobrania. Sieci z ćwiczenia 6 mają ich \textbf{ok. 3 miliony} (YOLO z listą 80 klas) i \textbf{ok. 14 milionów} (YOLOE, której zadajemy pytania słowami). Czy dałoby się je dobrać suwakami?
\item
  Do nauczenia takiej sieci potrzeba setek tysięcy opisanych zdjęć i wielu dni pracy kart graficznych. Dlaczego w praktyce korzysta się z \textbf{gotowych, wytrenowanych modeli}?
\end{enumerate}

\clearpage% Wygenerowano z ../06_detektor.md skryptem _md_do_latex.py (można edytować ręcznie).
\section{Ćwiczenie 6: gotowa sieć neuronowa szuka obiektów, czyli jak zapytać sieć (20 min)}\label{ux107wiczenie-6-gotowa-sieux107-neuronowa-szuka-obiektuxf3w-czyli-jak-zapytaux107-sieux107-20-min}

\textbf{Notatnik:} \texttt{06\_detektor.ipynb}

\textbf{Potrzebne:} obiekt do szukania (np. czerwona gumowa kaczka) na podłodze przed robotem, kilka innych przedmiotów (krzesło, plecak, pudełko), miarka, „obiekty-zagadki'' od prowadzącego

\subsection{Cel}\label{cel}

Sprawdzić, co potrafią gotowe sieci rozpoznające obiekty, nauczyć się dobierać zapytanie słowne i wybrać ustawienia do zadania „znajdź obiekt'' w ćwiczeniu 7.

\subsection{Trochę teorii}\label{trochux119-teorii}

Porównamy dwa rodzaje gotowych, wytrenowanych sieci:

\begin{itemize}
\tightlist
\item
  \textbf{Sieć z zamkniętą listą klas (YOLO, zbiór COCO).} Nauczona na ponad 100 000 opisanych zdjęć rozpoznaje \textbf{80 rodzajów obiektów}: osobę, krzesło, plecak, butelkę, kubek i inne. Czego nie ma na liście, tego nie wykryje. Najwyżej pomyli to z czymś podobnym.
\item
  \textbf{Sieć z otwartym słownikiem (YOLOE).} Tej sieci mówisz słowami (po angielsku), czego ma szukać, np. \texttt{"rubber\ duck"}. Nauczyła się łączyć obrazy z opisami słownymi, więc znajduje obiekty, na których \textbf{nikt jej specjalnie nie uczył}. Wynik zależy jednak od tego, \textbf{jak zapytasz}.
\end{itemize}

Dla każdego wykrytego obiektu sieć podaje \textbf{ramkę} (gdzie on jest) i \textbf{pewność} (od 0 do 1). Wykrycia słabsze od \textbf{progu pewności} są odrzucane.

Sieci działają na \textbf{procesorze} stacji. Jedno zdjęcie analizują przez kilkaset milisekund, czyli ok. 2--10 zdjęć na sekundę.

\subsection{Kroki}\label{kroki}

\subsubsection{Krok 1: sieć z zamkniętą listą (3 min)}\label{krok-1-sieux107-z-zamkniux119tux105-listux105-3-min}

Postaw kaczkę ok. 1 m przed robotem, a obok krzesło i plecak. Uruchom komórkę z modelem \texttt{"coco"}.

\begin{itemize}
\tightlist
\item
  Co wykryła sieć? \linia{14.0}
\item
  Czy wykryła kaczkę? Jeśli tak, to jako co? \linia{9.2}
\end{itemize}

\subsubsection{Krok 2: zapytaj słowami (5 min)}\label{krok-2-zapytaj-sux142owami-5-min}

Uruchom komórkę \texttt{porownaj\_slowa}. Sprawdza ona wiele opisów na tym samym zdjęciu i pokazuje, przy którym sieć jest najpewniejsza. Dopisz do listy \textbf{własne pomysły} (co najmniej 4).

\par\noindent\begin{minipage}{\linewidth}
\begin{tabular}{@{}
  >{\raggedright\arraybackslash}p{(\linewidth - 4\tabcolsep) * \real{0.3846}}
  >{\raggedright\arraybackslash}p{(\linewidth - 4\tabcolsep) * \real{0.3077}}
  >{\raggedright\arraybackslash}p{(\linewidth - 4\tabcolsep) * \real{0.3077}}@{}}
\toprule\noalign{}
\begin{minipage}[b]{\linewidth}\raggedright
Opis (zapytanie)
\end{minipage} & \begin{minipage}[b]{\linewidth}\raggedright
Pewność
\end{minipage} & \begin{minipage}[b]{\linewidth}\raggedright
Uwagi
\end{minipage} \\
\midrule\noalign{}
\texttt{duck} & & \\
\texttt{rubber\ duck} & & \\
\texttt{kaczka} (po polsku) & & \\
\texttt{bird} & & \\
& & \\
& & \\
& & \\
& & \\
\bottomrule\noalign{}
\end{tabular}
\end{minipage}\par\medskip

Najlepsze zapytanie: \linia{9.2}

\subsubsection{Krok 3: słowa konkurują (2 min)}\label{krok-3-sux142owa-konkurujux105-2-min}

Ustaw w \texttt{ustaw\_slowa} jednocześnie Twoje najlepsze zapytanie i inne słowa, np. \texttt{{[}"rubber\ duck",\ "box",\ "toy"{]}}. Czy pewność wykrycia kaczki się zmieniła? Czy pojawiły się nowe wykrycia? \linia{9.2}

\subsubsection{Krok 4: obiekt-zagadka (4 min)}\label{krok-4-obiekt-zagadka-4-min}

Prowadzący daje zespołowi przedmiot. Ustawcie go przed robotem i znajdźcie zapytanie, przy którym sieć wykrywa go z pewnością \textbf{powyżej 0,5}. Wygrywa zespół, który zrobi to najszybciej.

Przedmiot: \linia{6.7} najlepsze zapytanie: \linia{6.7} pewność: \linia{2.5}

\subsubsection{Krok 5: odległość, rozdzielczość, szybkość (4 min)}\label{krok-5-odlegux142oux15bux107-rozdzielczoux15bux107-szybkoux15bux107-4-min}

Ustawiaj kaczkę w różnych odległościach. Dla każdej odległości uruchom komórkę, która sprawdza rozdzielczości 320 i 640.

\par\noindent\begin{minipage}{\linewidth}
\begin{tabular}{@{}
  >{\raggedright\arraybackslash}p{(\linewidth - 4\tabcolsep) * \real{0.2549}}
  >{\raggedright\arraybackslash}p{(\linewidth - 4\tabcolsep) * \real{0.3725}}
  >{\raggedright\arraybackslash}p{(\linewidth - 4\tabcolsep) * \real{0.3725}}@{}}
\toprule\noalign{}
\begin{minipage}[b]{\linewidth}\raggedright
Odległość {[}m{]}
\end{minipage} & \begin{minipage}[b]{\linewidth}\raggedright
Pewność przy 320 px
\end{minipage} & \begin{minipage}[b]{\linewidth}\raggedright
Pewność przy 640 px
\end{minipage} \\
\midrule\noalign{}
0,5 & & \\
1,0 & & \\
2,0 & & \\
3,0 & & \\
\bottomrule\noalign{}
\end{tabular}
\end{minipage}\par\medskip

Szybkość: 320 px: \linia{2.5} zdjęć/s, 640 px: \linia{2.5} zdjęć/s

\subsubsection{Krok 6: decyzja zespołu (2 min)}\label{krok-6-decyzja-zespoux142u-2-min}

Wybierzcie ustawienia do ćwiczenia 7, w którym robot ma sam znaleźć obiekt i do niego podjechać.

\begin{itemize}
\tightlist
\item
  Zapytanie: \linia{6.7} próg: \linia{2.5} rozdzielczość: \linia{2.5}
\item
  Uzasadnienie: \linia{14.0}
\end{itemize}

\subsection{Pytania do przemyślenia}\label{pytania-do-przemyux15blenia}

\begin{enumerate}
\def\labelenumi{\arabic{enumi}.}
\tightlist
\item
  Dlaczego sieć z zamkniętą listą nie może wykryć kaczki, choć „widzi'' ją na zdjęciu?
\item
  Dlaczego zapytanie po polsku działa gorzej niż po angielsku? \emph{(Podpowiedź: na jakich opisach uczono sieć?)}
\item
  Robot ma szukać obiektu w całym pomieszczeniu. Co jest ważniejsze: duża rozdzielczość (widać daleko) czy szybkość? Czy da się połączyć jedno i drugie, np. szukać dokładnie, a podjeżdżać szybko?
\item
  Co jest gorsze dla robota szukającego obiektu: \textbf{fałszywe wykrycie} (pojedzie do czegoś innego) czy \textbf{przeoczenie}? Jak to wpływa na wybór progu?
\end{enumerate}

\clearpage% Wygenerowano z ../07_nawigacja.md skryptem _md_do_latex.py (można edytować ręcznie).
\section{Ćwiczenie 7: autonomiczna jazda na mapie, czyli zadania na torze (35 min)}\label{ux107wiczenie-7-autonomiczna-jazda-na-mapie-czyli-zadania-na-torze-35-min}

\textbf{Notatnik:} \texttt{07\_nawigacja.ipynb}

\textbf{Potrzebne:} Wasza mapa \texttt{zespolN\_najlepsza} (albo mapa prowadzącego \texttt{tor\_prowadzacy}), punkty \textbf{S, A, B} oznaczone na Waszym torze (współrzędne na tablicy), karton jako „nowa przeszkoda'', obiekt do znalezienia (kaczka) ukryty przez prowadzącego

\subsection{Cel}\label{cel}

Uruchomić nawigację autonomiczną: robot sam ustala, gdzie jest na mapie, planuje trasę do celu, omija przeszkody, a na koniec sam odnajduje obiekt w pomieszczeniu i do niego podjeżdża.

\subsection{Trochę teorii}\label{trochux119-teorii}

Nawigacja składa się z trzech elementów:

\begin{enumerate}
\def\labelenumi{\arabic{enumi}.}
\tightlist
\item
  \textbf{Lokalizacja (AMCL)}: robot ustala, gdzie jest na \textbf{gotowej} mapie, porównując skan LiDAR-u z mapą. Używa wielu hipotez położenia, czyli \textbf{cząstek} (na podglądzie mapy to chmura niebieskich strzałek). Na początku cząstki są rozrzucone. Gdy robot jedzie, te, które nie pasują do skanu, znikają, a chmura się zawęża.
\item
  \textbf{Planer globalny}: wyznacza trasę od położenia robota do celu po mapie (zielona linia na podglądzie mapy).
\item
  \textbf{Planer lokalny}: prowadzi robota wzdłuż trasy i na bieżąco omija przeszkody, które widzi LiDAR, także te, których nie było na mapie.
\end{enumerate}

Współrzędne mapy: punkt \textbf{S} to (0, 0), oś \textbf{x} biegnie wzdłuż linii startowej, oś \textbf{y} w lewo. Jednostką są metry.

\subsection{Kroki}\label{kroki}

\subsubsection{Krok 1: start nawigacji i lokalizacja (7 min)}\label{krok-1-start-nawigacji-i-lokalizacja-7-min}

\begin{enumerate}
\def\labelenumi{\arabic{enumi}.}
\tightlist
\item
  Uruchom \texttt{uruchom\_nawigacje("zespolN\_najlepsza")}. Start trwa do 1 minuty.
\item
  Postaw robota na punkcie \textbf{S} i uruchom \texttt{ustaw\_pozycje\_startowa(0,\ 0,\ 0)}.
\item
  Uruchom \texttt{panel\_sterowania(podglad=True)} i obejrzyj na mapie chmurę cząstek. Przejedź kawałek, obracając się powoli. Co się dzieje z chmurą?
\item
  Uruchom \texttt{gdzie\_jestem()}. Jaka jest niepewność położenia przed jazdą i po niej?
\end{enumerate}

Niepewność przed jazdą: \linia{2.5} cm, po jeździe: \linia{2.5} cm

\subsubsection{Zadanie A: dojazd do celu (7 min)}\label{zadanie-a-dojazd-do-celu-7-min}

Ustaw w komórce współrzędne punktu \textbf{A} (z tablicy) i uruchom \texttt{jedz\_do(x,\ y)}. Gdy robot stanie, zmierzcie miarką odległość środka robota od punktu A. Powtórz dla punktu \textbf{B}.

\subsubsection{Zadanie B: nowa przeszkoda (7 min)}\label{zadanie-b-nowa-przeszkoda-7-min}

\begin{enumerate}
\def\labelenumi{\arabic{enumi}.}
\tightlist
\item
  Postaw karton na trasie z A do B, w miejscu, gdzie na mapie jest pusto.
\item
  Wyślij robota do celu i obserwuj na podglądzie mapy (pojawia się pod komórką \texttt{jedz\_do}), jak planer zmienia trasę.
\item
  Zmień parametry komórką \texttt{ustaw\_parametry\_nawigacji(...)} i powtórz:

  \begin{itemize}
  \tightlist
  \item
    \texttt{odleglosc\_bezpieczenstwa}: 0,10 m, potem 0,35 m. Jak blisko przeszkód jedzie robot? Czy zmieści się w wąskim przejeździe?
  \item
    \texttt{maks\_predkosc}: 0,10 m/s, potem 0,25 m/s.
  \end{itemize}
\end{enumerate}

\subsubsection{Zadanie C: znajdź kaczkę (12 min)}\label{zadanie-c-znajdux17a-kaczkux119-12-min}

Prowadzący ukrywa kaczkę gdzieś na torze. Robot ma ją \textbf{sam} znaleźć i do niej podjechać.

Kaczka jest niższa niż wiązka LiDAR-u, więc LiDAR jej nie widzi, a nawigacja nie traktuje jej jak przeszkody. Robot szuka jej \textbf{kamerą i siecią z ćwiczenia 6}: obraca się, robi zdjęcia i sprawdza, czy sieć coś wykrywa. Gdy znajdzie kaczkę, podjeżdża: ustawia się tak, żeby kaczka była na środku obrazu, jedzie kawałek i powtarza. Zatrzymuje się, gdy kaczka jest już nisko w obrazie, czyli blisko robota.

\begin{enumerate}
\def\labelenumi{\arabic{enumi}.}
\tightlist
\item
  \textbf{C1: rozglądanie się (3 min).} Ustaw kaczkę ok. 1 m od robota, z boku albo z tyłu. Ustaw zapytanie i próg z ćwiczenia 6, a potem uruchom \texttt{rozejrzyj\_sie()}. Po ilu obrotach robot ją znalazł? Sprawdź też inny \texttt{krok} obrotu (np. 45° zamiast 30°).
\item
  \textbf{C2: dojazd (3 min).} Uruchom \texttt{podjedz\_do\_obiektu()}. Zmierz miarką, jak daleko od kaczki robot się zatrzymał. Zmień \texttt{dol\_stop}, żeby zatrzymał się ok. 15--20 cm przed kaczką.
\item
  \textbf{C3: poszukiwania (6 min).} Kaczka jest ukryta na torze. Na swojej mapie wybierzcie \textbf{punkty obserwacyjne}, z których robot się rozejrzy, i ustalcie ich kolejność. Uruchom \texttt{szukaj\_obiektu(punkty)}. Robot rozejrzy się na miejscu, a potem będzie jeździł po punktach, aż znajdzie kaczkę i do niej podjedzie. Komórka wypisze dziennik z czasem.
\end{enumerate}

\par\noindent\begin{minipage}{\linewidth}
\begin{tabular}{@{}
  >{\raggedright\arraybackslash}p{(\linewidth - 10\tabcolsep) * \real{0.0840}}
  >{\raggedright\arraybackslash}p{(\linewidth - 10\tabcolsep) * \real{0.1597}}
  >{\raggedright\arraybackslash}p{(\linewidth - 10\tabcolsep) * \real{0.1345}}
  >{\raggedright\arraybackslash}p{(\linewidth - 10\tabcolsep) * \real{0.1681}}
  >{\raggedright\arraybackslash}p{(\linewidth - 10\tabcolsep) * \real{0.2521}}
  >{\raggedright\arraybackslash}p{(\linewidth - 10\tabcolsep) * \real{0.2017}}@{}}
\toprule\noalign{}
\begin{minipage}[b]{\linewidth}\raggedright
Próba
\end{minipage} & \begin{minipage}[b]{\linewidth}\raggedright
Punkty obserwacyjne
\end{minipage} & \begin{minipage}[b]{\linewidth}\raggedright
krok {[}°{]}
\end{minipage} & \begin{minipage}[b]{\linewidth}\raggedright
Czas znalezienia {[}s{]}
\end{minipage} & \begin{minipage}[b]{\linewidth}\raggedright
Czas do zatrzymania przy kaczce {[}s{]}
\end{minipage} & \begin{minipage}[b]{\linewidth}\raggedright
Odległość od kaczki {[}cm{]}
\end{minipage} \\
\midrule\noalign{}
1 & & & & & \\
2 & & & & & \\
\bottomrule\noalign{}
\end{tabular}
\end{minipage}\par\medskip

Wygrywa zespół, którego robot najszybciej dojedzie do kaczki.

\subsubsection{Dziennik prób}\label{dziennik-pruxf3b}

Pracujcie w cyklu \textbf{test → analiza → korekta parametrów → kolejny test}.

\par\noindent\begin{minipage}{\linewidth}
\begin{tabular}{@{}
  >{\raggedright\arraybackslash}p{(\linewidth - 12\tabcolsep) * \real{0.0917}}
  >{\raggedright\arraybackslash}p{(\linewidth - 12\tabcolsep) * \real{0.0917}}
  >{\raggedright\arraybackslash}p{(\linewidth - 12\tabcolsep) * \real{0.1743}}
  >{\raggedright\arraybackslash}p{(\linewidth - 12\tabcolsep) * \real{0.1468}}
  >{\raggedright\arraybackslash}p{(\linewidth - 12\tabcolsep) * \real{0.1468}}
  >{\raggedright\arraybackslash}p{(\linewidth - 12\tabcolsep) * \real{0.2018}}
  >{\raggedright\arraybackslash}p{(\linewidth - 12\tabcolsep) * \real{0.1468}}@{}}
\toprule\noalign{}
\begin{minipage}[b]{\linewidth}\raggedright
Próba
\end{minipage} & \begin{minipage}[b]{\linewidth}\raggedright
Zadanie
\end{minipage} & \begin{minipage}[b]{\linewidth}\raggedright
Zmienione parametry
\end{minipage} & \begin{minipage}[b]{\linewidth}\raggedright
Dotarł? (T/N)
\end{minipage} & \begin{minipage}[b]{\linewidth}\raggedright
Czas {[}s{]}
\end{minipage} & \begin{minipage}[b]{\linewidth}\raggedright
Odległość od celu {[}cm{]}
\end{minipage} & \begin{minipage}[b]{\linewidth}\raggedright
Co się stało?
\end{minipage} \\
\midrule\noalign{}
1 & A & domyślne & & & & \\
2 & & & & & & \\
3 & & & & & & \\
4 & & & & & & \\
5 & & & & & & \\
6 & & & & & & \\
\bottomrule\noalign{}
\end{tabular}
\end{minipage}\par\medskip

\subsection{Pytania do przemyślenia}\label{pytania-do-przemyux15blenia}

\begin{enumerate}
\def\labelenumi{\arabic{enumi}.}
\tightlist
\item
  Czym różni się \textbf{lokalizacja} (AMCL) od \textbf{budowy mapy} (SLAM)? Dlaczego w tym ćwiczeniu nie uruchamialiśmy SLAM?
\item
  Co się stanie, gdy podasz robotowi \textbf{błędną} pozycję startową (np. 1 m obok)? Sprawdźcie, jeśli starczy czasu.
\item
  Dlaczego robot podjeżdża do kaczki „skokami'' (jedzie, staje, patrzy)? Wymieńcie wszystkie źródła opóźnienia (kamera → Wi-Fi → sieć neuronowa → decyzja → silniki).
\item
  Jak wybraliście punkty obserwacyjne? Co by się stało, gdyby kaczka stała za pudłem?
\item
  Kto odpowiada za to, gdy autonomiczny pojazd popełni błąd: programista, producent czujnika, właściciel czy ktoś inny?
\end{enumerate}

\begin{quote}
\textbf{Jeśli coś nie działa.}

\begin{itemize}
\tightlist
\item
  „PRZERWANO, nie da się dojechać'': cel leży w przeszkodzie albo zbyt blisko ściany. Zmniejsz \texttt{odleglosc\_bezpieczenstwa} albo wybierz inny punkt.
\item
  Robot kręci się w miejscu: lokalizacja się zgubiła. Ustaw ponownie pozycję startową.
\item
  Robot „nie widzi'' kaczki, choć stoi ona przed nim: obniż próg, zmień zapytanie albo ustaw rozdzielczość 640.
\item
  Robot jedzie do czegoś innego: podnieś próg albo zmień zapytanie na dokładniejsze.
\item
  Awaryjnie: \texttt{zatrzymaj\_wszystko()}.
\end{itemize}
\end{quote}

\clearpage% Wygenerowano z ../08_karta_wynikow.md skryptem _md_do_latex.py (można edytować ręcznie).
\section{Ćwiczenie 8: karta wyników zespołu (20 min, prezentacja ok. 2 min na zespół)}\label{ux107wiczenie-8-karta-wynikuxf3w-zespoux142u-20-min-prezentacja-ok.-2-min-na-zespuxf3ux142}

\textbf{Zespół nr:} \linia{1.7} \textbf{Skład:} \linia{14.0}

\subsection{1. Czujniki}\label{czujniki}

\begin{itemize}
\tightlist
\item
  Największy błąd LiDAR-u, który zmierzyliśmy: \linia{2.5} cm, przy \linia{9.2}
\item
  Materiał, przy którym LiDAR zawiódł: \linia{9.2}
\item
  Błąd odometrii po jednym kwadracie 1 × 1 m: \linia{2.5} cm, po dwóch: \linia{2.5} cm
\end{itemize}

\subsection{2. Mapa (SLAM)}\label{mapa-slam}

\begin{itemize}
\tightlist
\item
  Najlepsze parametry: rozdzielczość \linia{1.7} zasięg \linia{1.7} aktualizacja \linia{1.7} cząstki \linia{1.7}
\item
  Co najbardziej psuło mapę: \linia{14.0}
\end{itemize}

\szkicmapy

\subsection{3. Sieć neuronowa i zapytania}\label{sieux107-neuronowa-i-zapytania}

\begin{itemize}
\tightlist
\item
  Najlepsze zapytanie dla kaczki: \linia{6.7} pewność: \linia{2.5}
\item
  Obiekt-zagadka: \linia{6.7} zapytanie: \linia{6.7} pewność: \linia{2.5}
\item
  Wybrane ustawienia: próg \linia{2.5} rozdzielczość \linia{2.5} zdjęć/s: \linia{2.5}
\item
  Kiedy sieć się myliła: \linia{14.0}
\end{itemize}

\subsection{4. Autonomiczna jazda}\label{autonomiczna-jazda}

\par\noindent\begin{minipage}{\linewidth}
\begin{tabular}{@{}
  >{\raggedright\arraybackslash}p{(\linewidth - 6\tabcolsep) * \real{0.1915}}
  >{\raggedright\arraybackslash}p{(\linewidth - 6\tabcolsep) * \real{0.1702}}
  >{\raggedright\arraybackslash}p{(\linewidth - 6\tabcolsep) * \real{0.3191}}
  >{\raggedright\arraybackslash}p{(\linewidth - 6\tabcolsep) * \real{0.3191}}@{}}
\toprule\noalign{}
\begin{minipage}[b]{\linewidth}\raggedright
Zadanie
\end{minipage} & \begin{minipage}[b]{\linewidth}\raggedright
Udało się?
\end{minipage} & \begin{minipage}[b]{\linewidth}\raggedright
Najlepszy wynik (czas / dokładność)
\end{minipage} & \begin{minipage}[b]{\linewidth}\raggedright
Najważniejsza zmiana parametrów
\end{minipage} \\
\midrule\noalign{}
A: dojazd do celu & & & \\
B: nowa przeszkoda & & & \\
C: znajdź kaczkę & & & \\
\bottomrule\noalign{}
\end{tabular}
\end{minipage}\par\medskip

\subsection{5. Wnioski (do prezentacji)}\label{wnioski-do-prezentacji}

\begin{itemize}
\tightlist
\item
  Największy problem i jego przyczyna: \linia{14.0}
\item
  Co poprawilibyśmy, mając więcej czasu: \linia{14.0}
\item
  Gdzie w prawdziwym świecie widzicie zastosowanie tego, co robiliśmy: \linia{10.1}
\end{itemize}


\end{document}
